\begin{table}[htbp]
\centering
\caption{Robustness: Employment Effects Under Alternative Specifications}
\label{tab:robust}
\begin{tabular}{lcccc}
\toprule
Specification & Coefficient & SE & $p$-value & $N$ \\
\midrule
\textit{Baseline (border counties)} & $-$0.117*** & (0.014) & $<$0.001 & 19,726 \\
Pre-COVID (2015--2019) & $-$0.045** & (0.016) & 0.034 & 7,901 \\
Trimmed window (drop 2017--2018) & $-$0.094** & (0.031) & 0.026 & 13,812 \\
Sector $\times$ Quarter FE & $-$0.118*** & (0.017) & 0.001 & 19,726 \\
Leave-one-state-out range & [$-$0.128, $-$0.108] & --- & --- & --- \\
Placebo (2017Q1) & 0.065** & (0.020) & 0.022 & 5,914 \\
Simple DiD (no industry) & $-$0.050 & (0.033) & 0.192 & 19,726 \\
RI $p$-value (state permutation) & --- & --- & 0.167 & --- \\
RI $p$-value (timing permutation) & --- & --- & 0.077 & --- \\
\bottomrule
\end{tabular}
\begin{tablenotes}[flushleft]
\small
\item \textit{Notes:} $^{***}p<0.01$, $^{**}p<0.05$, $^{*}p<0.10$. All specifications use log employment as the dependent variable with standard errors clustered at the state level. Baseline: county-sector and quarter FE ($N = 19{,}726$). Pre-COVID restricts to 2015--2019. Trimmed window drops 2017--2018. Sector $\times$ Quarter FE absorbs nationwide sector-specific trends. LOSO drops each of five control states in turn. RI permutes treatment state (5 placebos) and timing (pre-period quarters).\end{tablenotes}
\end{table}
